
\vspace{-5pt}
\section{Related Work}
\label{sec:related_work}
\vspace{-5pt}

\textbf{Poison Detection in LLM Code Generation.}
There are multiple work investigating the malicious behavior 
of LLMs in the inference stage for code generation. 
The work in~\citet{zeng2025inducing} studies a 
poisoning attack in code generation when the
external tools including search engines used by LLMs 
contain malicious content. 
Similarly, BIPIA~\citep{yi2025benchmarking} presents 
the first systematic benchmark to evaluate 
indirect prompt injection attacks, focusing 
on malicious instructions embedded in external 
content that manipulate LLM behavior. 
In contrast, our work demonstrates a more fundamental 
problem that does not require the LLM to 
access any external sources during inference. 
We show that harmful content, such as scam API endpoints, 
has already been absorbed into the models' weights from their 
training data and can be triggered by developer-style prompts.


\textbf{Poisoning Attacks in LLM Training Pipelines.} 
Data poisoning, where adversaries manipulate 
training data to alter model behavior at 
inference, has emerged as a critical threat to 
machine learning systems. While early work mostly 
focused on computer vision 
applications~\citep{cina2023wild, raghavan2022improved, 
goldblum2022dataset}, 
recent studies have extended this concern to the 
language domain, particularly LLMs. 
\citet{carlini2024poisoning} 
proves the practicality and feasibility of 
poisoning web-scale training datasets collection pipelines.
A recent survey~\citep{zhao2025data} provides the 
first systematic overview of data poisoning attacks 
targeting LLMs across multiple stages such as 
pretraining, fine-tuning, preference alignment, and 
instruction tuning. \citet{jiang2024turning} studies 
poisoning attacks on LLMs 
during fine-tuning on text summarization and completion tasks, 
showing that existing defenses remain ineffective. 
Our study differs by (i) targeting malicious code 
generation, which poses immediate execution risks, rather 
than natural-language outputs, and (ii) auditing production 
LLMs for evidence of existing, passive poisoning in 
their training corpora, rather than new inference-time active attacks.


\chen{\textbf{Evaluating LLM-generated Code Security from CWEs}}
There is another line of work that constructs test benchmarks 
from existing Common Weakness Enumeration (CWE) entries to 
evaluate the security of LLM-generated code. 
\textsc{SecCodePLT}~\citep{yang2024seccodeplt} utilizes seed 
samples of real-world vulnerable code and employs LLM-based mutators 
for data expansion across 44 CWEs. 
Similarly, \textsc{SafeGenBench}~\citep{li2025safegenbench}, 
based on pre-defined vulnerability categories and CWE types, 
applies LLMs to generate test questions that are not only consistent 
with real development scenarios but also strictly adhering 
to specific vulnerability characteristics. 
In a slightly different direction, the \textsc{CodeLMSec} 
benchmark~\citep{debenedetti2023a} approximates an inversion 
of the target black-box model via few-shot prompting to 
build non-secure prompts for evaluation. 
Unlike these benchmarks, which focus on specific programming 
languages and CWEs, our work investigates a distinct, orthogonal threat: 
the generation of malicious URLs. We target the model's 
memorization and reproduction of real-world scams, rather than 
flaws in code logic. 




% \textbf{Uncurated Datasets.}
% Deep learning models, particularly LLMs, achieve their 
% best performance when trained on massive datasets, 
% as demonstrated by neural scaling 
% laws~\citep{kaplan2020scaling, hoffmann2022training, 
% brown2020language, carlini2024poisoning}. 
% While some of the most advanced language models, 
% such as GPT-4 and Gemini Ultra, do not 
% disclose the size of their training datasets, 
% estimates based on their training compute suggest 
% they were likely trained on approximately 13 trillion 
% tokens—assuming Chinchilla scaling 
% efficiency~\citep{hoffmann2022training, 
% villalobos2024position}. As a result, the 
% demand for publicly available human-written 
% text is growing rapidly. Recent projections 
% suggest that if current development trends 
% continue, LLMs may exhaust the available 
% supply of public human text between 2026 and 2032, 
% or even earlier under continued 
% over-training~\citep{villalobos2024position}. 
% To meet these growing data demands, 
% researchers increasingly turn to large-scale 
% web scraping to expand their training corpora, 
% raising new concerns about data quality, and security.

% \textbf{Large Language Models for code generation.} 
% Large language models have demonstrated 
% remarkable capabilities in generating code, 
% enabling them to assist in software 
% development tasks such as code completion, 
% bug fixing, and even synthesizing entire 
% programs from natural language prompts. 
% Recent models trained explicitly on 
% extensive code corpora, such as 
% OpenAI's Codex\citep{chen2021evaluating}, 
% DeepSeek-Coder\citep{zhu2024deepseek, 
% guo2024deepseek}, and Google's Gemini\citep{team2023gemini}, 
% show that domain-targeted training enhances 
% performance on Python-focused code 
% generation benchmarks like 
% HumanEval\citep{chen2021evaluating} and 
% MBPP\citep{austin2021program}. Claude\citep{TheC3} from 
% Anthropic has also been evaluated for its 
% reasoning and code generation abilities, 
% especially in multi-turn interactions. These 
% advances have spurred a wave of research into 
% both the capabilities and limitations of 
% code-generating LLMs, as well as the potential 
% security risks posed by automatically generated 
% code, especially in cases where the training of a 
% model is outsourced\citep{cina2023wild}.

